\scanpage{99}
\mspage{\handwritten{4}}
\hfill\unsure{20.2.}\\
I finally went on pondering about a definition of $\infty$-groupoids, and \handwritten{it} seems to me that, after all, the topological motivation does furnish the ``simple guiding principle'' which yesterday seemed to me to be still to be discovered, \inserted{\handwritten{in order}} not to get lost in the \unsure{messyness|messiness} of ever higher order structures. Let me try to put it down roughly.

\handwritten{4.}\quad First I would like to correct somewhat the rather undiscriminate description I gave yesterday of what I thought of viewing as ``primary'', secondary, ternary etc structures for an $\infty$-groupoid. More careful refletion conduces to view as the most primitive, starting structure on the set of sets $F_i$ ($i\in\underline{N}$), as a skeleton on which progressively organs and flesh will be added, the mere data of the source and target maps
\[
s_i,t_i:F_i\rightrightarrows F_{i-1}\qquad(i\geq1)\ ,
\]
which it will be convenient to supplement \deleted{by the} formally by corresponding maps $s_0,t_0$ for $i=1$, from $F_0$ to $F_{-1}\mathrel{\overset{\mathrm{def}}{=}}$ one-element set. In a moment we will pass to a universal situation, when the $F_i$ are replaced by the corresponding ``universal'' objects $\underline{F}_i$ in a suitable catégory stable under finite products, where $F_{-1}$ will be the final element. For several reasons, it is not proper to view the inversion maps $\mathrm{inv}_i:F_i\to F_i$, and still less the other $i-1$ involutions on $F_i$ which I at first overlooked\handwritten{,} as being part of the primitive or ``skeletal'' structure. One main reason is that already for the most usual $2$-groupoids, such as the $2$-groupoid whose $0$-objects are ordinary ($1$-) groupoids, the $1$-objects being equivalences between these (namely functors which are fully faithful and essentially \unsure{surjectve|surjective}), and the $2$-objects morphisms (or ``natural transformations'') between such, there is \uline{not}, for an $1$-object $f:C\to C'$, a natural choice of an ``inverse'' namely of a quasi-inverse in the usual sense. And even assuming that such quasi-inverse is chosen for every $f$, it is by no means clear that such choice can be made involutive, namely such that $(f^{\vee})^{\vee}=f$ for every $f$ (and not merely $(f^{\vee})^{\vee}$ isomorphic to $f$). The maps $\mathrm{inv}_i$ will appear rather, quite naturally, as ``primary structure'', and they will not be involutions, but ``pseudo-involutions'' (namely involution ``up to homotopy''). It turns out that among the various functors \inserted{\handwritten{$\Pi_\infty$}} we will construct\handwritten{,} from the category of topological spaces to the category of $\infty$-groupoids\deleted{\unsure{,}} \handwritten{(}the construction depending on arbitrary choices and \unsure{yielding|yilding} a large bunch of mutually non isomorphic functors, which however are ``equivalent'' in a sense which will have to be made precise\unsure{\ldots}\deleted{-}\handwritten{)} there are choices neater than others, and some of these will yield in the primary structure maps $\mathrm{inv}_i$ which are actual involutions, and similarly for the other pseudo-involutions, appearing \deleted{\unsure{as}}\inserted{in succession as} higher order structure. The possibility of such neat and fairly natural choices had somewhat mislead me
% archive sheet number: 98
